This present work addresses the creation of \textit{RoadVet}, a mobile application aimed at home-care veterinary professionals to register medical records, tutors, and pets. The motivation to create this application is explained by the network inconsistency problems present in remote areas, these problems compromise the agility and security of clinical care for pets, as well as cause the loss of important data for veterinary professionals. To develop this solution, the \textit{offline-first} paradigm was used, which provides an adequate structure for the system, allowing data to be saved on the client-side in cases of connection failure. Methodologically, the study is characterized as applied research with experimental development. The solution's architecture was divided between the client-side layer, implemented with the \textit{React Native} library, along with \textit{WatermelonDB} for local data persistence, and the support layer, implemented using the \textit{Django Rest Framework} and the \textit{SQLite} relational database. The asynchronous synchronization process was designed to operate through an event queue, adopting the Last-Write-Wins conflict resolution strategy. Empirical tests carried out in a controlled environment proved that the system is capable of maintaining complete user operability without connectivity and of synchronizing data in an integral and automated way after network reestablishment. It is concluded that the \textit{offline-first} approach mitigates the risks of data loss and increases the efficiency of mobile veterinary care.

% Separe as Keywords por ponto
\keywords{Mobile application. \textit{Offline-first}. \textit{React Native}. \textit{WatermelonDB}. Data synchronization.}